\documentclass[10pt,aspectratio=169]{beamer}

\usepackage[UTF8]{ctex}
\useinnertheme{metropolis}
\useoutertheme{metropolis}
\usecolortheme{metropolis}
\usefonttheme{professionalfonts}
\usepackage{appendixnumberbeamer}
\usepackage{booktabs}
\usepackage{amsmath,amssymb,bm}
\usepackage{graphicx}
\usepackage{array}
\usepackage{tabularx}
\usepackage{ragged2e}
\usepackage{xcolor}

\title{Deep-RKPM-KAN}
\subtitle{从 1D 机制验证到 2D Poisson 收敛}
\author{陈一锋}
\institute{Deepritz Enhanced}
\date{2026-03-08}

\setsansfont{Segoe UI}
\setmonofont{Consolas}
\setCJKmainfont{Microsoft YaHei UI}
\setCJKsansfont{Microsoft YaHei UI}

\newcommand{\smallnote}[1]{\vspace{0.3em}{\scriptsize\color{gray}#1}}
\newcommand{\imgw}[2]{\includegraphics[width=#1\linewidth]{#2}}
\newcommand{\imgfull}[1]{\includegraphics[width=0.92\linewidth]{#1}}
\newcommand{\takeaway}[1]{\textbf{结论：}#1}

\begin{document}

\maketitle

\begin{frame}{目录}
  \setbeamertemplate{section in toc}[sections numbered]
  \tableofcontents[hideallsubsections]
\end{frame}

\section{背景与目标}

\begin{frame}{研究目标：学习可解释的 meshfree shape function}
\begin{columns}[T,onlytextwidth]
\column{0.54\textwidth}
\begin{block}{核心问题}
传统 RKPM/MLS 强结构但依赖显式局部矩阵；黑箱神经 PDE 方法优化方便但弱化了离散空间解释性。
\end{block}
\begin{itemize}
  \item 我们希望直接学习具有无网格结构约束的形函数
  \item 目标不仅是求解 PDE，还要保留 \texttt{PU}、一次重构与边界行为
  \item 训练路径分为形函数预训练与 Deep Ritz 物理阶段
\end{itemize}
\column{0.42\textwidth}
\begin{alertblock}{关键词}
\begin{itemize}
  \item Partition of unity
  \item Linear reproduction
  \item Boundary anchoring
  \item Teacher distillation
  \item Stability / robustness
\end{itemize}
\end{alertblock}
\end{columns}
\smallnote{叙事主线：1D 机制验证 $\rightarrow$ 2D patch test $\rightarrow$ 2D 主实验 $\rightarrow$ 稳定性与鲁棒性。}
\end{frame}

\begin{frame}{方法总图：从节点到 Deep Ritz}
\begin{center}
\renewcommand{\arraystretch}{1.35}
\begin{tabular}{>{\RaggedRight\arraybackslash}p{0.16\linewidth}>{\centering\arraybackslash}p{0.05\linewidth}>{\RaggedRight\arraybackslash}p{0.16\linewidth}>{\centering\arraybackslash}p{0.05\linewidth}>{\RaggedRight\arraybackslash}p{0.16\linewidth}>{\centering\arraybackslash}p{0.05\linewidth}>{\RaggedRight\arraybackslash}p{0.22\linewidth}}
节点集 & $\rightarrow$ & 局部 KAN 母核 & $\rightarrow$ & 窗函数与归一化 & $\rightarrow$ & 形函数 $\Phi_I(x)$ 与 Deep Ritz \\
\end{tabular}
\end{center}
\vspace{0.55em}
{\footnotesize
\begin{align*}
\tilde\phi_I(x) &= \psi_\theta\!\left(\frac{x-x_I}{\kappa h}\right) W\!\left(\|x-x_I\|,\kappa h\right),\\
\phi_I(x) &= \frac{\tilde\phi_I(x)}{\sum_J \tilde\phi_J(x)},
\qquad
u_h(x)=\sum_I \phi_I(x) w_I.
\end{align*}
}
\begin{itemize}
  \item Phase A：先学习 shape function 的几何一致性与边界锚定
  \item Phase B：在保持离散空间稳定的前提下最小化 Deep Ritz 能量
  \item fallback、teacher、regularization 是关键稳定化机制
\end{itemize}
\smallnote{$\kappa h$ 控制支撑半径，$W$ 为紧支撑窗函数，$w_I$ 为节点系数。}
\end{frame}

\section{1D 机制验证}

\begin{frame}{为什么先做 1D？}
\begin{columns}[T,onlytextwidth]
\column{0.55\textwidth}
\begin{itemize}
  \item 1D 更适合直接观察 shape function 的形状、边界行为与振荡
  \item 可以清楚分辨 \texttt{softplus}、边界锚定、teacher 的作用
  \item 是进入 2D 前的“显微镜”试验
\end{itemize}
\column{0.4\textwidth}
\begin{exampleblock}{本节回答的问题}
\begin{itemize}
  \item 什么情况下 shape function 会失稳？
  \item 什么机制能恢复稳定？
  \item 哪种设计值得推广到 2D？
\end{itemize}
\end{exampleblock}
\end{columns}
\vspace{0.4em}
\takeaway{1D 是机制验证平台，不只是预备实验。}
\end{frame}

\begin{frame}{1D baseline：基本 shape quality 已经成立}
\begin{columns}[T,onlytextwidth]
\column{0.47\textwidth}
\imgw{1.0}{../../output/meshfree_kan_rkpm_1d_validation/deepritz_meshfree_kan_rkpm_1d_validation_output/shape_compare.png}
\column{0.47\textwidth}
\imgw{1.0}{../../output/meshfree_kan_rkpm_1d_validation/deepritz_meshfree_kan_rkpm_1d_validation_output/shape_subset_overlay.png}
\end{columns}
\vspace{0.3em}
\begin{itemize}
  \item \texttt{global\_l2}=2.26e-2
  \item \texttt{pu\_max\_error}=1.47e-12
  \item \texttt{linear\_reproduction\_rmse}=2.56e-3
\end{itemize}
\smallnote{数据来源：\texttt{output/meshfree\_kan\_rkpm\_1d\_validation/.../metrics.json}。}
\end{frame}

\begin{frame}{1D failure case：去掉 softplus 后会明显失稳}
\begin{columns}[T,onlytextwidth]
\column{0.47\textwidth}
\imgw{1.0}{../../output/meshfree_kan_rkpm_1d_validation/deepritz_meshfree_kan_rkpm_1d_validation_v4_output_no_softplus_v4/shape_compare.png}
\column{0.47\textwidth}
\imgw{1.0}{../../output/meshfree_kan_rkpm_1d_validation/deepritz_meshfree_kan_rkpm_1d_validation_v4_output_no_softplus_v4/shape_subset_overlay.png}
\end{columns}
\vspace{0.3em}
\begin{itemize}
  \item \texttt{global\_l2}=3.435
  \item \texttt{global\_linf}=19.295
  \item \texttt{pu\_raw\_sum\_rmse}=1.033
\end{itemize}
\takeaway{单纯移除稳定化先验会导致形函数质量灾难性退化。}
\end{frame}

\begin{frame}{1D teacher 版本：结构恢复并更适合推广到 2D}
\begin{columns}[T,onlytextwidth]
\column{0.47\textwidth}
\imgw{1.0}{../../output/meshfree_kan_rkpm_1d_validation/deepritz_meshfree_kan_rkpm_1d_validation_v5_output_teacher_distill_v5/shape_compare.png}
\column{0.47\textwidth}
\imgw{1.0}{../../output/meshfree_kan_rkpm_1d_validation/deepritz_meshfree_kan_rkpm_1d_validation_v5_output_teacher_distill_v5/shape_subset_overlay.png}
\end{columns}
\vspace{0.3em}
\begin{itemize}
  \item \texttt{global\_l2}=4.91e-2
  \item \texttt{linear\_reproduction\_rmse}=7.50e-3
  \item $\phi_0(0)\approx0.961,\; \phi_N(1)\approx0.954$
\end{itemize}
\takeaway{teacher 约束为 shape-function 学习提供了更可控的优化路径。}
\end{frame}

\begin{frame}{1D 小结：2D 设计应保留哪些机制？}
\begin{enumerate}
  \item 形函数学习需要结构化约束，而不是完全自由拟合
  \item teacher、fallback、regularization 是 2D 稳定性的关键候选
  \item 因此 2D 不只比较 PDE 误差，还要报告 \texttt{PU}、$\Lambda_h$ 与一次重构
\end{enumerate}
\vspace{0.5em}
\begin{alertblock}{本节结论}
1D 实验给出的不是最终结果，而是 2D 方案设计的依据。
\end{alertblock}
\end{frame}

\section{2D 形函数一致性与主实验}

\begin{frame}{2D patch test：先验证线性重构一致性}
\begin{columns}[T,onlytextwidth]
\column{0.32\textwidth}
\imgw{1.0}{../../output/meshfree_kan_rkpm_2d_validation/patch_test/variant_softplus_raw_pu_bd_ns15_k2p5_seed42_paper_v1/figures/pu_heatmap.png}
\column{0.32\textwidth}
\imgw{1.0}{../../output/meshfree_kan_rkpm_2d_validation/patch_test/variant_softplus_raw_pu_bd_ns15_k2p5_seed42_paper_v1/figures/linear_x_heatmap.png}
\column{0.32\textwidth}
\imgw{1.0}{../../output/meshfree_kan_rkpm_2d_validation/patch_test/variant_softplus_raw_pu_bd_ns15_k2p5_seed42_paper_v1/figures/lambda_h_heatmap.png}
\end{columns}
\vspace{0.15em}
{\scriptsize
\begin{align*}
\text{PU residual} &: \sum_I \phi_I(x,y)-1,\qquad
\text{linear reproduction}: \sum_I \phi_I(x,y)x_I-x,\\
\Lambda_h(x,y) &: \sum_I |\phi_I(x,y)|.
\end{align*}
}
\begin{itemize}
  \item 前两图检验 PU / 一次重构，第三图检验是否存在放大效应
  \item $\Lambda_h\approx 1$ 常对应稳定形函数；若显著偏大，通常说明有振荡或抵消
\end{itemize}
\end{frame}

\begin{frame}{2D patch test：为什么这一步很关键？}
\begin{columns}[T,onlytextwidth]
\column{0.58\textwidth}
\begin{align*}
\sum_I \phi_I(x,y) &\approx 1,\\
\sum_I \phi_I(x,y)x_I &\approx x,\\
\sum_I \phi_I(x,y)y_I &\approx y.
\end{align*}
\column{0.38\textwidth}
\begin{exampleblock}{意义}
\begin{itemize}
  \item 常数重构成立
  \item 一次 patch 重构成立
  \item 说明进入 PDE 阶段前，离散空间已具备基本一致性
\end{itemize}
\end{exampleblock}
\end{columns}
\vspace{0.5em}
\takeaway{patch test 让 2D 主实验具备更强的可解释性。}
\end{frame}

\begin{frame}{2D Poisson 主实验设定}
\begin{columns}[T,onlytextwidth]
\column{0.56\textwidth}
{\footnotesize
\begin{align*}
u^*(x,y) &= \sin(\pi x)\sin(\pi y),\\
-\Delta u &= f = 2\pi^2 \sin(\pi x)\sin(\pi y),\\
\mathcal{L}_{\mathrm{PDE}} &= \mathcal{E}(u_h)+\beta_{bc}\,\mathcal{L}_{bc}+\gamma_{lin}\,\mathcal{L}_{lin}.
\end{align*}
}
\begin{itemize}
  \item 网格规模：\texttt{n\_side = 11, 15, 21}
  \item 当前主变体：\texttt{no\_softplus\_teacher\_reg}
  \item \texttt{kappa=2.5}：支撑半径比，$r_{supp}=\kappa h$
  \item \texttt{beta\_bc=500}：边界罚项权重
  \item \texttt{gamma\_linear=1}：一次重构保持项权重
\end{itemize}
\column{0.38\textwidth}
\begin{block}{两阶段训练}
\begin{itemize}
  \item Phase A：先把 $\phi_I$ 训练到 PU / 边界 / teacher 一致性较好
  \item Phase B：在解 PDE 的同时继续抑制离散空间漂移
\end{itemize}
\end{block}
\begin{exampleblock}{目标}
让 2D \texttt{L2/H1} 结果进入更适合投稿展示的误差区间，同时保持 \texttt{PU} 与 $\Lambda_h$ 稳定。
\end{exampleblock}
\end{columns}
\smallnote{当前最佳结果对应 \texttt{paper\_v3\_main}；这里更强调参数含义与损失组成，而不只报最终误差。}
\end{frame}

\begin{frame}{2D 主结果：多尺度误差对比}
\begin{columns}[T,onlytextwidth]
\column{0.60\textwidth}
\imgw{1.0}{../../output/meshfree_kan_rkpm_2d_validation/poisson_convergence/convergence_summary.png}
\column{0.36\textwidth}
\begin{itemize}
  \item \texttt{n=11}: \texttt{L2}=7.19e-3\\ \texttt{H1}=2.77e-1
  \item \texttt{n=15}: \texttt{L2}=8.53e-3\\ \texttt{H1}=2.36e-1
  \item \texttt{n=21}: \texttt{L2}=1.15e-2\\ \texttt{H1}=2.73e-1
\end{itemize}
\begin{alertblock}{要点}
\small 这页展示的是不同离散尺度下的当前最佳误差水平；目前还不宜表述为严格单调的 $h$-收敛。
\end{alertblock}
\end{columns}
\end{frame}

\begin{frame}{2D 主结果：同尺度下旧版到最优版的改进}
\centering
\small
\begin{tabular}{@{}lcc@{}}
\toprule
尺度 & baseline \texttt{paper\_v1} & best \texttt{paper\_v3\_main} \\
\midrule
\texttt{n=11} & \texttt{L2: 2.58e-2}, \texttt{H1: 5.58e-1} & \texttt{L2: 7.19e-3}, \texttt{H1: 2.77e-1} \\
\texttt{n=15} & \texttt{L2: 3.02e-2}, \texttt{H1: 5.40e-1} & \texttt{L2: 8.53e-3}, \texttt{H1: 2.36e-1} \\
\texttt{n=21} & \texttt{L2: 2.73e-2}, \texttt{H1: 6.18e-1} & \texttt{L2: 1.15e-2}, \texttt{H1: 2.73e-1} \\
\bottomrule
\end{tabular}
\vspace{0.7em}
\begin{alertblock}{关键信息}
同尺度对比更能说明调参收益：当前改进主要体现在误差区间明显下降，而不是已经获得漂亮的跨尺度单调收敛。
\end{alertblock}
\end{frame}

\begin{frame}{2D 主结果：代表性解场可视化}
\begin{columns}[T,onlytextwidth]
\column{0.64\textwidth}
\imgw{1.0}{../../output/meshfree_kan_rkpm_2d_validation/poisson_convergence/variant_no_softplus_teacher_reg_ns15_k2p5_seed42_paper_v3_main/figures/solution_triplet.png}
\column{0.32\textwidth}
{\footnotesize \textbf{Phase A：shape function 预训练}}\\[-0.2em]
\imgw{1.0}{../../output/meshfree_kan_rkpm_2d_validation/poisson_convergence/variant_no_softplus_teacher_reg_ns15_k2p5_seed42_paper_v3_main/figures/loss_curves_phase_a.png}
\vspace{0.35em}

{\footnotesize \textbf{Phase B：Poisson 能量优化}}\\[-0.2em]
\imgw{1.0}{../../output/meshfree_kan_rkpm_2d_validation/poisson_convergence/variant_no_softplus_teacher_reg_ns15_k2p5_seed42_paper_v3_main/figures/loss_curves_phase_b.png}
\end{columns}
\vspace{0.4em}
\takeaway{双图展示后，Phase A 与 Phase B 的优化目标被清晰分离，更适合答辩讲解。}
\end{frame}
\section{稳定性与鲁棒性}

\begin{frame}{稳定性消融：teacher 与 regularization 为什么重要？}
\begin{columns}[T,onlytextwidth]
\column{0.60\textwidth}
\imgw{1.0}{../../output/meshfree_kan_rkpm_2d_validation/stability_ablation/stability_summary.png}
\column{0.36\textwidth}
\begin{itemize}
  \item \texttt{softplus\_raw\_pu\_bd}: $\Lambda_h^{\max}\approx1.00$
  \item \texttt{no\_softplus\_raw\_pu\_bd}: $\Lambda_h^{\max}\approx14.09$
  \item \texttt{no\_softplus\_teacher\_reg}: $\Lambda_h^{\max}\approx1.01$
\end{itemize}
\begin{alertblock}{要点}
\small teacher guidance 恢复了 raw unconstrained 版本失去的稳定性。
\end{alertblock}
\end{columns}
\end{frame}

\begin{frame}{稳定性消融：结论如何解释？}
\begin{columns}[T,onlytextwidth]
\column{0.53\textwidth}
\begin{itemize}
  \item 不加约束地去掉 softplus，会引入明显的放大与抵消
  \item fallback 能缓解局部覆盖退化，但 teacher 更直接约束了 shape matrix
  \item regularization 进一步抑制原始响应过大造成的病态
\end{itemize}
\column{0.43\textwidth}
\begin{exampleblock}{最重要的信号}
$\Lambda_h$ 从 14.09 回落到接近 1，说明改进不是偶然拟合，而是结构层面的稳定恢复。
\end{exampleblock}
\end{columns}
\end{frame}

\begin{frame}{不规则节点鲁棒性：轻度扰动下依然可用}
\begin{columns}[T,onlytextwidth]
\column{0.60\textwidth}
\imgw{1.0}{../../output/meshfree_kan_rkpm_2d_validation/irregular_nodes/irregular_summary.png}
\column{0.36\textwidth}
\begin{itemize}
  \item \texttt{jitter=0.0}: \texttt{L2}=2.41e-2\\ \texttt{H1}=6.64e-1
  \item \texttt{jitter=0.1}: \texttt{L2}=2.65e-2\\ \texttt{H1}=6.26e-1
  \item \texttt{jitter=0.2}: \texttt{L2}=2.63e-2\\ \texttt{H1}=7.58e-1
\end{itemize}
\begin{alertblock}{要点}
\small 中等节点扰动下，误差上升有限，说明方法具备一定鲁棒性。
\end{alertblock}
\end{columns}
\end{frame}

\begin{frame}{不规则节点：代表性解场}
\begin{columns}[T,onlytextwidth]
\column{0.64\textwidth}
\imgw{1.0}{../../output/meshfree_kan_rkpm_2d_validation/irregular_nodes/variant_softplus_raw_pu_bd_ns15_k2p5_jit0p1_seed42_paper_v1/figures/solution_triplet.png}
\column{0.32\textwidth}
{\footnotesize \textbf{Phase A：shape function 预训练}}\\[-0.2em]
\imgw{1.0}{../../output/meshfree_kan_rkpm_2d_validation/irregular_nodes/variant_softplus_raw_pu_bd_ns15_k2p5_jit0p1_seed42_paper_v1/figures/loss_curves_phase_a.png}
\vspace{0.35em}

{\footnotesize \textbf{Phase B：Poisson 能量优化}}\\[-0.2em]
\imgw{1.0}{../../output/meshfree_kan_rkpm_2d_validation/irregular_nodes/variant_softplus_raw_pu_bd_ns15_k2p5_jit0p1_seed42_paper_v1/figures/loss_curves_phase_b.png}
\end{columns}
\vspace{0.4em}
\takeaway{双图训练轨迹更容易说明：即使在轻度节点扰动下，两阶段优化依然稳定。}
\end{frame}
\section{总结}

\begin{frame}{总结}
\begin{enumerate}
  \item 1D 实验明确了 shape-function 学习所需的关键稳定化机制
  \item 2D patch test 说明离散空间在进入 PDE 前已具备 PU 与一次重构一致性
  \item 2D \texttt{no\_softplus\_teacher\_reg} 是当前最强投稿候选配置
\end{enumerate}
\vspace{0.5em}
\begin{alertblock}{当前最强信息}
这个方法的优势不仅体现在误差更低，还体现在：稳定、可解释、可做系统消融。
\end{alertblock}
\end{frame}

\begin{frame}[standout]
谢谢！\\[0.5em]
答疑时可继续展示 appendix 中的更多参数扫描与对比图。
\end{frame}

\appendix

\begin{frame}{Backup：可继续补充的内容}
\begin{itemize}
  \item 更多 \texttt{paper\_v1 / paper\_v2 / paper\_v3} 参数演化表
  \item \texttt{n=11,15,21} 三组 \texttt{solution\_triplet} 横向对照
  \item 1D 更多版本的 \texttt{shape\_compare} 与边界锚定指标
  \item 与 RKPM / FEM / PINN 的进一步横向比较
\end{itemize}
\end{frame}

\end{document}





